\documentclass[10pt, aspectratio=169]{beamer}

\usepackage{fontspec}
\usepackage[russian]{babel}

\setsansfont{Arial} 
\setmonofont{Courier New} 

\usepackage{xcolor}
\usepackage{listings}
\usepackage{booktabs}
\usepackage{multicol}

\definecolor{mainblue}{RGB}{40, 50, 180}
\definecolor{codebg}{RGB}{245, 245, 245}
\definecolor{codegray}{RGB}{120, 120, 120}
\definecolor{stringred}{RGB}{180, 40, 40}
\definecolor{keywordblue}{RGB}{0, 0, 255}

\setbeamercolor{palette primary}{bg=mainblue, fg=white}
\setbeamercolor{title}{bg=mainblue, fg=white}
\setbeamercolor{frametitle}{bg=mainblue, fg=white}
\setbeamercolor{structure}{fg=mainblue}
\setbeamercolor{itemize item}{fg=mainblue}

\setbeamertemplate{footline}{
  \begin{beamercolorbox}[wd=\paperwidth, ht=3ex, dp=1.5ex, leftskip=3mm, rightskip=3mm]{palette primary}
    \fontsize{6pt}{6pt}\selectfont Максим Емельянов, Литвинов Юрий \hfill С++ Семинар \hfill \insertframenumber{} / \inserttotalframenumber
  \end{beamercolorbox}
}
\setbeamertemplate{navigation symbols}{} 

\lstset{
    backgroundcolor=\color{codebg},
    basicstyle=\ttfamily\scriptsize,
    keywordstyle=\color{keywordblue}\bfseries,
    stringstyle=\color{stringred},
    commentstyle=\color{green!50!black},
    numbers=none,
    breaklines=true,
    frame=single,
    rulecolor=\color{gray!30},
    showstringspaces=false,
    tabsize=4,
    keepspaces=true
}

\title{Динамический полиморфизм в C++}
\subtitle{Механика vtable, цена абстракции и современные альтернативы}
\author{Максим Емельянов \and Литвинов Юрий}
\date{C++ Семинар \\ \today}

\begin{document}

\begin{frame}[plain]
    \titlepage
\end{frame}

\begin{frame}{План}
    \tableofcontents
\end{frame}

\section{Часть 1: Механика vtable}

\begin{frame}{Что такое полиморфизм? Раннее vs Позднее связывание}
    \begin{itemize}
        \item \textbf{Динамический полиморфизм}~--- механизм, при котором программа определяет точный адрес вызываемой функции во время выполнения программы (Runtime), исходя из реального типа объекта в памяти.
        \vspace{0.3cm}
        \item \textbf{Раннее связывание (Early binding):}\\
        Вызов функции привязывается к указателю на этапе компиляции. Не подходит для восходящего преобразования (upcast).
        \vspace{0.3cm}
        \item \textbf{Позднее связывание (Late binding):}\\
        Гарантирует вызов правильной версии функции для конкретного объекта, независимо от типа ссылки. В С++ реализуется через \texttt{virtual}.
    \end{itemize}
\end{frame}

\begin{frame}[fragile]{Взгляд под капот: vptr и vtable}
    При добавлении \texttt{virtual} компилятор неявно добавляет в объект указатель \textbf{vptr}, который ссылается на \textbf{vtable} (таблицу виртуальных методов).

    \vspace{0.3cm}
    \textbf{Как выглядит вызов в ассемблере (GCC):}
\begin{lstlisting}[language={[x86masm]Assembler}]
mov rbp, QWORD PTR [rbx]   ; В регистр rbp помещаем указатель на объект
mov rax, QWORD PTR [rbp+0] ; В rax помещаем указатель на vtable - разыменование 1
call [QWORD PTR [rax]]     ; Вызов функции по адресу из vtable - разыменование 2
\end{lstlisting}

    \begin{itemize}
        \item \textbf{Преимущество:} Правильный вызов метода дочернего класса.
        \item \textbf{Недостаток:} Затраты на двойное обращение к памяти.
    \end{itemize}
\end{frame}

\section{Часть 2: Классический подход и шаблоны}

\begin{frame}[fragile]{Классический пример: Иерархия (С++98)}
\begin{lstlisting}[language=C++]
class Property {
protected:
    double worth;
public:
    Property(double worth) : worth(worth) {}
    virtual double getTax() const = 0; // Чисто виртуальный метод
};

class Car : public Property {
public:
    Car(double worth) : Property(worth) {}
    double getTax() const override { return worth / 200; }
};

// Использование:
Property* p = new Car(400000);
std::cout << p->getTax(); // Вызовется Car::getTax()
\end{lstlisting}
\end{frame}

\begin{frame}[fragile]{Статическая имитация (Шаблоны)}
Можно избежать \texttt{vtable}, используя шаблоны:
\begin{lstlisting}[language=C++]
template <typename T>
class Base {
public:
    void someFunctionality() {
        static_cast<T*>(this)->someFunctionality();
    }
};

class Child : public Base<Child> {
public:
    void someFunctionality() {
        cout << "Child class perform new function..." << endl;
    }
};
\end{lstlisting}
    \textbf{Итог:} Ускоряет работу (нет vtable), но может сильно раздуть бинарник (каждый тип генерирует новый класс).
\end{frame}

\section{Часть 3: std::variant}

\begin{frame}[fragile]{Альтернатива C++17: std::variant + std::visit}
\begin{lstlisting}[language=C++]
class Derived {
public: void PrintName() const { std::cout << "calling Derived!\n"; }
};
class ExtraDerived {
public: void PrintName() const { std::cout << "calling ExtraDerived!\n"; }
};

int main() {
    std::variant<Derived, ExtraDerived> var = Derived{};
    
    // Вызов через лямбду (visitor)
    auto caller = [](const auto& obj) { obj.PrintName(); };
    std::visit(caller, var);
}
\end{lstlisting}
\end{frame}

\begin{frame}{Плюсы и минусы std::variant}
    \begin{columns}[T]
        \begin{column}{0.5\textwidth}
            \textbf{\color{mainblue}Преимущества:}
            \begin{itemize}
                \item \textbf{Value semantics:} объекты хранятся без \texttt{new/delete}.
                \item \textbf{Нет базы:} классы могут быть не связаны иерархией.
                \item \textbf{Duck typing:} сигнатуры не обязаны строго совпадать.
            \end{itemize}
        \end{column}
        
        \begin{column}{0.5\textwidth}
            \textbf{\color{stringred}Недостатки:}
            \begin{itemize}
                \item Все типы должны быть известны \textbf{заранее} (плохо для плагинов).
                \item \textbf{Перерасход памяти:} размер равен самому большому типу + тег (теряем байты на мелких объектах).
            \end{itemize}
        \end{column}
    \end{columns}
\end{frame}
\section{Часть 4: Полиморфизм в Go}

\begin{frame}[fragile]{Альтернативный взгляд: Полиморфизм в Go}
    \begin{itemize}
        \item \textbf{Duck Typing}: неявная реализация интерфейсов.
        \item Нет наследования.
        \item Если тип содержит методы интерфейса~--- он его реализует.
    \end{itemize}

\begin{lstlisting}[language=Go, basicstyle=\tiny, backgroundcolor=\color{gray!10}]
type Geometry interface {
    Area() float64
}

type Rect struct { Width, Height float64 }
func (r Rect) Area() float64 { return r.Width * r.Height }

func Measure(g Geometry) {
    fmt.Println(g.Area())
}

func main() {
    r := Rect{Width: 3, Height: 4}
    Measure(r) // Работает автоматически
}
\end{lstlisting}
\end{frame}

\begin{frame}[fragile]{Взгляд на реализацию интерфейсов в Go}
    Переменная интерфейса в Go (\texttt{iface})~--- это структура из двух указателей:
    
\begin{lstlisting}[language=Go, basicstyle=\small, backgroundcolor=\color{gray!10}]
type iface struct {
    tab  *itab          // Указатель на таблицу интерфейсов
    data unsafe.Pointer // Указатель на реальные данные
}
\end{lstlisting}

    \begin{itemize}
        \item \textbf{itab (Interface Table)}: содержит метаданные типа и список указателей на функции. Схожа с vtable в C++.
        \item \textbf{Динамика и кэш}: itab генерируется в рантайме при первом сопоставлении пары "Интерфейс-Тип" и кэшируется.
        \item \textbf{Экономия}: в отличие от C++, объект не содержит \texttt{vptr}, если он не используется через интерфейс.
    \end{itemize}
\end{frame}

\section{Часть 5: Полиморфизм в Rust}


\begin{frame}[fragile]
\frametitle{Динамический полиморфизм в Rust: Traits}
\begin{itemize}
    \item \textbf{Трейты (Traits):} Описывают общее поведение типов.
    \item \textbf{dyn:} Ключевое слово для обозначения динамической диспетчеризации.
    \item \textbf{Box<dyn Trait>:} Так как размер объекта неизвестен при компиляции, работа всегда идет через указатель.
\end{itemize}

\begin{lstlisting}[language=C++, basicstyle=\tiny, keywordstyle=\color{blue}]
trait Animal { fn make_sound(&self); }
struct Dog;
impl Animal for Dog {
    fn make_sound(&self) { println!("Woof!"); }
}

fn main() {
    let animals: Vec<Box<dyn Animal>> = vec![
        Box::new(Dog),
        Box::new(Cat),
    ];
    for a in animals { a.make_sound(); }
}
\end{lstlisting}
\end{frame}


\begin{frame}
\frametitle{Под капотом Rust: Fat Pointers}
В отличие от C++, Rust не хранит \texttt{vptr} внутри объекта. Вместо этого используется \textbf{толстый указатель} (16 байт на x64):

\vspace{0.5cm}
\begin{center}
\begin{tabular}{|c|c|}
\hline
\textbf{Data Pointer (8 byte)} & \textbf{vtable Pointer (8 byte)} \\ \hline
Адрес данных в памяти & Адрес таблицы методов \\ \hline
\end{tabular}
\end{center}
\vspace{0.5cm}

\begin{itemize}
    \item \textbf{Преимущество:} Объекты «чистые», не занимают лишнюю память, если не используются полиморфно.
    \item \textbf{vtable содержит:}
    \begin{itemize}
        \item Указатель на деструктор.
        \item Размер типа и выравнивание (\textit{size, align}).
        \item Адреса реализаций методов трейта.
    \end{itemize}
\end{itemize}
\end{frame}

\begin{frame}
\frametitle{Rust dyn против Go Interface}
Хотя оба языка имеют похожие механизмы полиморфизма, имеются следующие различия:
\vspace{0.3cm}
\begin{enumerate}
    \item \textbf{Статика (Rust) vs Динамика (Go):}
    \begin{itemize}
        \item В Rust vtable строится \textbf{на этапе компиляции}. Вызов максимально дешев.
        \item В Go \texttt{itab} может генерироваться \textbf{в рантайме} (через хеш-таблицы) при первом приведении типов.
    \end{itemize}
    \item \textbf{Явность (Explicit) vs Неявность (Implicit):}
    \begin{itemize}
        \item В Go реализация неявная: совпали сигнатуры~--- интерфейс реализован.
        \item В Rust требуется явный \texttt{impl Trait for Struct}, что исключает ошибки «случайного» совпадения имен.
    \end{itemize}
\end{enumerate}
\end{frame}

\section{Часть 6: Сравнение подходов}

\begin{frame}
\frametitle{Итоговое сравнение подходов}
\small
\begin{table}
\centering
\begin{tabular}{|l|l|l|l|l|}
\hline
\textbf{Критерий} & \textbf{C++ Virtual} & \textbf{C++ variant} & \textbf{Go} & \textbf{Rust} \\ \hline
\textbf{Связывание} & Runtime (vtable) & Runtime (index) & Runtime (itab) & Runtime (vtable) \\ \hline
\textbf{Разметка} & Содержит vptr & Данные + тег & Чистый объект & Чистый объект \\ \hline
\textbf{Расширение} & Легко & Сложно & Легко & Легко \\ \hline
\textbf{Инлайнинг} & Почти никогда & Часто & Редко & Редко \\ \hline
\end{tabular}
\end{table}
\end{frame}

\begin{frame}
\frametitle{Заключение: Что быстрее?}
\begin{itemize}
    \item \textbf{std::variant}~--- потенциально самый быстрый за счет локальности данных и инлайнинга (для небольшого набора типов).
    
    \item \textbf{vtable (C++)}~--- стандарт для больших иерархий. Основные затраты: хранение \texttt{vptr} в каждом объекте и двойное разыменовывание при вызове.
    
    \item \textbf{Rust dyn}~--- предлагает баланс: объекты остаются «чистыми» (без \texttt{vptr}), а таблицы методов генерируются во время компиляции. Это дает производительность уровня C++.
    
    \item \textbf{Интерфейсы Go}~--- гибкий подход (Duck Typing). Удобен для крупных систем, но несет небольшие расходы на построение таблиц \texttt{itab} в рантайме.

    \item \textbf{Выбор зависит от задачи:} 
    \begin{itemize}
        \item Максимальная скорость: \texttt{std::variant} / Rust.
        \item Гибкость архитектуры: Go / Rust.
        \item Сложные иерархии наследования: C++ (\texttt{virtual}).
    \end{itemize}
\end{itemize}
\end{frame}

\end{document}